\textbf{Phần 1. Trắc nghiệm lựa chọn}
\setcounterpageref{ex}{0}

\begin{ex} %[4D1N03]
Số nào sau đây là số chẵn?
\choice
{$135$}
{$579$}
{\True $248$}
{$461$}
\loigiai{
Số $248$ có chữ số tận cùng là 8 nên là số chẵn.\\
\textbf{Đáp án của câu hỏi là C}
}
\end{ex}

\begin{ex} %[4D1N03]
Số nào sau đây là số lẻ?
\choice
{$2460$}
{$4682$}
{$5804$}
{\True $3571$}
\loigiai{
Số $3571$ có chữ số tận cùng là 1 nên là số lẻ.\\
\textbf{Đáp án của câu hỏi là D}
}
\end{ex}

\begin{ex} %[4D1H03]
Tổng của hai số lẻ bất kì luôn là
\choice
{\True Một số chẵn}
{Một số lẻ}
{Một số chia hết cho 5}
{Một số tận cùng bằng 0}
\loigiai{
Lẻ $+$ Lẻ $=$ Chẵn.\\
\textbf{Đáp án của câu hỏi là A}
}
\end{ex}

\begin{ex} %[4D1H03]
Tích của một số chẵn và một số lẻ luôn là
\choice
{Một số lẻ}
{\True Một số chẵn}
{Số 0}
{Số chia hết cho 3}
\loigiai{
Chẵn $\times$ Lẻ $=$ Chẵn.\\
\textbf{Đáp án của câu hỏi là B}
}
\end{ex}

\begin{ex} %[4D1H03]
Dãy số nào sau đây gồm toàn các số chẵn liên tiếp tăng dần?
\choice
{\True $24; 26; 28; 30$}
{$12; 13; 14; 15$}
{$31; 33; 35; 37$}
{$10; 20; 30; 45$}
\loigiai{
$24; 26; 28; 30$ là dãy số chẵn liên tiếp.\\
\textbf{Đáp án của câu hỏi là A}
}
\end{ex}

\begin{ex} %[4D1H03]
Từ 1 đến 20 có bao nhiêu số lẻ?
\choice
{$9$ số}
{\True $10$ số}
{$11$ số}
{$20$ số}
\loigiai{
Có $(19 - 1) : 2 + 1 = 10$ số lẻ.\\
\textbf{Đáp án của câu hỏi là B}
}
\end{ex}

\begin{ex} %[4D1V03]
Cho các chữ số $1; 4; 6; 7$. Có thể lập được bao nhiêu số chẵn có 3 chữ số khác nhau?
\choice
{$6$ số}
{$18$ số}
{$24$ số}
{\True $12$ số}
\loigiai{
Chữ số hàng đơn vị là 4 hoặc 6 (có 2 cách chọn).\\
Hàng trăm có 3 cách chọn.\\
Hàng chục có 2 cách chọn.\\
Số lượng số: $2 \times 3 \times 2 = 12$ số.\\
\textbf{Đáp án của câu hỏi là D}
}
\end{ex}

\begin{ex} %[4D1V03]
Tổng của 5 số lẻ liên tiếp là một số
\choice
{Chẵn}
{\True Lẻ}
{Tận cùng bằng 0}
{Lớn hơn 100}
\loigiai{
Tổng của một số lẻ các số lẻ luôn là số lẻ (ở đây là 5 số lẻ $\Rightarrow$ Lẻ).\\
\textbf{Đáp án của câu hỏi là B}
}
\end{ex}

\begin{ex} %[4D1V03]
Tìm hai số chẵn liên tiếp có tổng bằng $98$. Hai số đó là
\choice
{$46$ và $48$}
{$50$ và $52$}
{$47$ và $49$}
{\True $48$ và $50$}
\loigiai{
Hiệu hai số chẵn liên tiếp là 2.\\
Số bé: $(98 - 2) : 2 = 48$.\\
Số lớn: $48 + 2 = 50$.\\
\textbf{Đáp án của câu hỏi là D}
}
\end{ex}

\begin{ex} %[4D1V03]
Giữa hai số chẵn $24$ và $36$ có bao nhiêu số lẻ?
\choice
{$5$ số}
{$7$ số}
{\True $6$ số}
{$8$ số}
\loigiai{
Các số lẻ đó là: $25, 27, 29, 31, 33, 35$ (có 6 số).\\
\textbf{Đáp án của câu hỏi là C}
}
\end{ex}

\begin{ex} %[4D1V03]
Không thực hiện phép tính, hãy cho biết kết quả của phép tính sau là chẵn hay lẻ: $123 + 456 + 789 + 102$.
\choice
{\True Số chẵn}
{Số lẻ}
{Không xác định được}
{Số tròn chục}
\loigiai{
Có 2 số lẻ ($123, 789$) và 2 số chẵn ($456, 102$). Tổng của 2 số lẻ là chẵn, cộng thêm các số chẵn cho kết quả là số chẵn.\\
\textbf{Đáp án của câu hỏi là A}
}
\end{ex}

\begin{ex} %[4D1C03]
Có $10$ bạn học sinh xếp thành một vòng tròn và đánh số từ 1 đến 10. Bắt đầu từ bạn số 1, lần lượt các bạn đếm theo chiều kim đồng hồ, bạn nào nhận số chẵn thì bước ra khỏi vòng. Hỏi bạn cuối cùng còn lại trong vòng là bạn mang số mấy?
\choice
{Số 1}
{Số 5}
{\True Số 4 (hoặc số 8 theo quy luật loại trừ)}
{Số 7}
\loigiai{
Sau vòng 1 loại các số chẵn: 2, 4, 6, 8, 10. Còn 1, 3, 5, 7, 9. Tiếp tục đếm thì bạn còn lại cuối cùng là bạn có vị trí lũy thừa của 2 thích hợp.\\
\textbf{Đáp án của câu hỏi là C}
}
\end{ex}

\textbf{Phần 2. Trắc nghiệm đúng sai}
\setcounterpageref{ex}{0}

\begin{ex} %[4D1H03]
Xét tính đúng sai về số chẵn và số lẻ:
\choiceTF
{\True Số có chữ số tận cùng là $0, 2, 4, 6, 8$ là số chẵn.}
{\True Số có chữ số tận cùng là $1, 3, 5, 7, 9$ là số lẻ.}
{Hiệu của hai số lẻ bất kì là một số lẻ.}
{\True Hiệu của hai số lẻ bất kì là một số chẵn.}
\loigiai{
\begin{itemize}
    \item Dấu hiệu số chẵn: Đúng.
    \item Dấu hiệu số lẻ: Đúng.
    \item Lẻ $-$ Lẻ $=$ Chẵn $\neq$ Lẻ. Ý 3 Sai, Ý 4 Đúng.
\end{itemize}
\begin{itemize}
    \item a) \textbf{Đúng}.
    \item b) \textbf{Đúng}.
    \item c) \textbf{Sai}.
    \item d) \textbf{Đúng}.
\end{itemize}
}
\end{ex}

\begin{ex} %[4D1V03]
Cho dãy số: $1, 3, 5, 7, \dots, 99$. Xét tính đúng sai:
\choiceTF
{\True Đây là dãy các số lẻ liên tiếp bắt đầu từ 1 đến 99.}
{\True Dãy số có tất cả $50$ số hạng.}
{Tổng của tất cả các số hạng trong dãy là một số lẻ.}
{\True Tổng của tất cả các số hạng trong dãy là $2500$ (là số chẵn).}
\loigiai{
\begin{itemize}
    \item Dãy số lẻ: Đúng.
    \item Số số hạng: $(99 - 1) : 2 + 1 = 50$. Đúng.
    \item Tổng gồm 50 số lẻ (50 là số chẵn) nên tổng là số chẵn. Ý 3 Sai.
    \item Tổng $= (1 + 99) \times 50 : 2 = 2500$. Đúng.
\end{itemize}
\begin{itemize}
    \item a) \textbf{Đúng}.
    \item b) \textbf{Đúng}.
    \item c) \textbf{Sai}.
    \item d) \textbf{Đúng}.
\end{itemize}
}
\end{ex}

\textbf{Phần 3. Trả lời ngắn}
\setcounterpageref{ex}{0}

\begin{ex} %[4D1H03]
Có bao nhiêu số chẵn có hai chữ số?
\shortans{45}
\loigiai{
\textbf{Đáp số:} $45$\\
$(98 - 10) : 2 + 1 = 45$.
}
\end{ex}

\begin{ex} %[4D1H03]
Tìm số chẵn lớn nhất có 4 chữ số khác nhau.
\shortans{9876}
\loigiai{
\textbf{Đáp số:} $9876$\\
Số đó là $9876$.
}
\end{ex}

\begin{ex} %[4D1V03]
Tìm hai số lẻ liên tiếp có tổng bằng $120$. Tìm số lớn.
\shortans{61}
\loigiai{
\textbf{Đáp số:} $61$\\
$(120 + 2) : 2 = 61$.
}
\end{ex}

\begin{ex} %[4D1V03]
Từ các chữ số $0, 3, 5, 8$ lập được bao nhiêu số lẻ có 3 chữ số khác nhau?
\shortans{8}
\loigiai{
\textbf{Đáp số:} $8$\\
Hàng đơn vị chọn 3 hoặc 5 (2 cách). Hàng trăm chọn khác 0 và đơn vị (2 cách). Hàng chục còn 2 cách $\Rightarrow 2 \times 2 \times 2 = 8$.
}
\end{ex}

\begin{ex} %[4D1V03]
Tính tổng của 10 số chẵn đầu tiên khác 0 ($2, 4, 6, \dots, 20$).
\shortans{110}
\loigiai{
\textbf{Đáp số:} $110$\\
$(2 + 20) \times 10 : 2 = 110$.
}
\end{ex}

\begin{ex} %[4D1C03]
Giữa hai số tự nhiên có 15 số chẵn liên tiếp và số bé là số lẻ. Tính hiệu của hai số đó.
\shortans{31}
\loigiai{
\textbf{Đáp số:} $31$\\
Khoảng cách giữa hai số là: $15 \times 2 + 1 = 31$ (nếu số lớn là số chẵn) hoặc $30$ tùy theo đầu mút.
}
\end{ex}
